\section{Tests Beyond the Benchmark}
\label{sec:beyond}
% SWEVO expansion: generated supplementary tables placed in Sections 8 and 9 (\SuppTable{<label>}). The capture is global;
% the figure of mpce_supp_diag.tex (fig:D-psodyn) is discarded here, so only its tables are stored.
\begingroup\RenewEnviron{figure}[1][]{}\RenewEnviron{figure*}[1][]{}%
\endgroup% [flattened: removed \CaptureSuppTables{analysis/mpce_supp_diag.tex}]
% [flattened: removed \CaptureSuppTables{analysis/mpce_supp_direction.tex}]
% Captions of the direction-resolution tables refer to main-text tables as M-<label> (supplement build); in the main
% text these names are aliases of the labels themselves (resolved from the .aux of the previous run).
\makeatletter
\def\sw@mainalias#1{\@ifundefined{r@M-#1}{\@ifundefined{r@#1}{}{%
  \global\expandafter\let\csname r@M-#1\expandafter\endcsname\csname r@#1\endcsname}}{}}
\sw@mainalias{tab:friedman68}\sw@mainalias{tab:hr-site}
\makeatother
\subsection{Budget Scaling}
\label{sec:budget}
% >>>>> begin analysis/mpce_tab_feasbudget.tex
%% generated by mpce_results.py -- do not edit by hand
\begin{table}[!t]
\centering
\caption{Initialization and budget on the six largest benchmark cases (30 seeds).}
\label{tab:feasbudget}
\scriptsize\setlength{\tabcolsep}{2.4pt}
\begin{tabular}{lccccccc}
\toprule
& \multicolumn{2}{c}{Random init.} & \multicolumn{2}{c}{Feasible init.} & \multicolumn{3}{c}{Avg.\ rank} \\
\cmidrule(lr){2-3}\cmidrule(lr){4-5}\cmidrule(lr){6-8}
Method & Loss & Feas. & Loss & Feas. & 6,030 & 30,030 & 120,030 \\
\midrule
PSO-VNS & 4.084 & 100.0 & 3.938 & 100.0 & 1.50 & 1.17 & 2.17 \\
SSA-VNS & 4.678 & 95.0 & 3.955 & 100.0 & 4.33 & 4.17 & 4.17 \\
LX-SSA-VNS & 4.802 & 98.9 & 3.952 & 100.0 & 5.17 & 5.67 & 5.83 \\
RS-VNS & 4.830 & 93.3 & n/r & n/r & 4.67 & 6.33 & 4.67 \\
VNS & 4.950 & 100.0 & 3.810 & 100.0 & 6.50 & 4.83 & 3.17 \\
PSO & 4.181 & 100.0 & 5.549 & 100.0 & 1.83 & 2.50 & 5.83 \\
SSA & 4.828 & 80.0 & 4.947 & 100.0 & 8.33 & 8.67 & 8.83 \\
LX-SSA & 5.379 & 77.8 & 5.048 & 100.0 & 8.67 & 9.83 & 10.00 \\
DE & -- & 1.1 & 5.549 & 100.0 & 10.00 & 5.17 & 3.33 \\
MS-SLSQP & 4.522 & 99.4 & 4.294 & 100.0 & 4.00 & 6.67 & 7.00 \\
\bottomrule
\end{tabular}
\par\vspace{2pt}\parbox{\columnwidth}{\scriptsize Loss: mean wake loss (\%) of the feasible runs; Feas.: feasible runs (\%); both at 6,030 evaluations with random and feasibility-preserving initialization. Avg.\ rank at 6,030, 30,030 and 120,030 evaluations, random initialization. n/r: not run.}
\end{table}
% <<<<< end analysis/mpce_tab_feasbudget.tex

On the six largest benchmark cases (random starts), PSO-VNS keeps the best average rank at all three budgets (6,030, 30,030 and 120,030 evaluations: 1.50, 1.17 and 2.17; Table~\ref{tab:feasbudget}; Fig.~\ref{fig:budget}) and ranks first in 4, 5 and 3 of the six cases (rank differences not tested); PSO's mean wake loss minus that of PSO-VNS is +0.10, +0.18 and +0.27~pp (Table~\ref{S-tab:budget-cases}). The order of the other methods depends on the budget: from 6,030 to 120,030 evaluations, PSO falls back from an average rank of 1.83 to 5.83, behind SSA-VNS (4.17), and VNS moves up from 6.50 to 3.17; at 120,030, VNS, SSA-VNS, RS-VNS, LX-SSA-VNS and DE come within 0.17~pp of the mean wake loss of PSO-VNS (3.14\%).
% CHECK-FINAL [C26] [C29] [C34]: PSO-VNS best average rank at all three budgets; first-place counts generated per budget
% CHECK-FINAL [C37]: at 120,030 SSA-VNS ranks ahead of PSO; VNS, SSA-VNS, RS-VNS, LX-SSA-VNS and DE within 0.2 pp (\NBudCloseGapOneTwentyK) of PSO-VNS. PSO and VNS ranks are generated macros (not tested by C37).

\begin{figure}[!htbp]
\centering
\pipefig{0.64\textwidth}{figures_mpce/budget_scaling.pdf}{6cm}
\caption{Budget scaling: mean wake loss of the six largest benchmark cases and mean AEP loss of the Horns Rev~1 block (feasible runs; methods with at least half of their runs feasible) versus the budget (random initialization; 10 seeds for Horns Rev~1 at 120,030 evaluations).}
\label{fig:budget}
\end{figure}

The methods differ mainly in how much they gain from a larger budget (Fig.~\ref{fig:budget}): from 6,030 to 120,030 evaluations, the mean wake loss of VNS falls from 4.95\% to 3.17\%, that of PSO only from 4.18\% to 3.42\% (PSO-VNS: 4.08\% to 3.14\%). On the Horns Rev~1 block (last panel), the mean AEP loss of PSO-VNS also falls (7.07\%, 6.52\% and 6.40\%), and at 120,030 evaluations DE has the lowest.
% verified (manual): budget losses from mpce_numbers.tex (\NBudLoss...): VNS - PSO-VNS mean wake loss 0.87 / 0.40 / 0.03 pp at 6,030 / 30,030 / 120,030 (decreasing); drop 6,030 -> 120,030: VNS 1.78, PSO 0.76, PSO-VNS 0.94 pp; DE feasible share at 6,030 from Table feasbudget (1.1 %); DE plotted only at 30,030 and 120,030 (at least half of the runs feasible)
% CHECK-FINAL [C25]: Horns Rev 16: 30,030 and 120,030 mean AEP loss of PSO-VNS below 6,030; CHECK-DIR [F29]: DE highest mean AEP (lowest loss) at 120,030

\subsection{Feasibility-Preserving Initialization}
\label{sec:feasinit}
With feasibility-preserving initialization, the share of feasible runs rises for SSA-VNS (95.0\% to 100.0\%), SSA (80.0\% to 100.0\%), LX-SSA (77.8\% to 100.0\%) and DE (1.1\% to 100.0\%). VNS alone is then best in mean wake loss and average rank, and PSO-VNS ranks second (rank 2.83; loss 3.94\%, 4.08\% with random starts; Table~\ref{tab:feasbudget}; per case: Table~\ref{S-tab:feasinit-cases}). PSO and DE never improve on the best initial layout (210 runs each; Table~\ref{S-tab:D-feasstart}): their moves change all or most coordinates at once, so in our runs no trial is both feasible and better and the global best never moves (Lemma~\ref{lem:feas-improve}, Proposition~\ref{prop:feasible-start}); VNS, which moves one coordinate at a time, is not affected. The advantage of PSO-VNS thus depends on the initialization: from random starts, PSO may help mainly by reaching feasibility.
% CHECK-FINAL [C28] [C38]: feasible starts: VNS alone lowest loss and best rank; PSO-VNS second in rank; PSO and DE never move from the best initial layout
% CHECK-DIAG [D13]: PSO and DE never improve on the best initial layout (instrumented and all stored runs)
% CHECK-DIAG [D14]: all PSO candidates except the G re-evaluations and all \NDFsCandDE{} DE trials infeasible, none accepted, no pbest update; first-move median min spacing \NDFsFirstMinSpPSO{} m
% CHECK-DIAG [D16]: the feasible initial layouts are distinct (matched RMS distance to G > 50 m)

What follows from the algorithms can be separated from what is only observed. Let $L(\mathbf X)$ be the wake loss, the first term of~\eqref{eq:penalty}, so that $F_p=L+\mathrm{Pen}$ with $\mathrm{Pen}\ge0$.
\begin{lemma}[Improvement needs a feasible, better trial]
\label{lem:feas-improve}
If the global best $\mathbf G$ is feasible, a trial $\mathbf X$ replaces it only if $L(\mathbf X)<L(\mathbf G)$ and every constraint violation $g$ of $\mathbf X$ satisfies $(1+\mu g)^2<L(\mathbf G)$; in the benchmark cases, $g<4.6\times10^{-8}$ (m or m$^2$), far below the reporting tolerance of $10^{-6}$~m.
\end{lemma}
\begin{proof}
Replacement needs $F_p(\mathbf X)<F_p(\mathbf G)=L(\mathbf G)$, and $F_p(\mathbf X)$ is the sum of $L(\mathbf X)\ge0$ and the violation terms $(1+\mu g)^2>0$. In the benchmark cases, $L(\mathbf G)\le NE_{\rm ideal}+\varepsilon$ with $N\le15$, $E_{\rm ideal}\le14{,}045.74$ ($\varepsilon$: the small negative part of the power curve above cut-in), so $\sqrt{L(\mathbf G)}<459.1$ and $g<(\sqrt{L(\mathbf G)}-1)/\mu<4.6\times10^{-8}$ for $\mu=10^{10}$.
\end{proof}
% verified: (sqrt(15 x 14045.74) - 1)/1e10 = 4.58e-8 (as optA/supp_theory.tex, Lemma S-feas-improve; the small negative part of the power curve above cut-in, epsilon, does not change the bound)
\begin{proposition}[PSO from feasible starts]
\label{prop:feasible-start}
Let all initial layouts be feasible, let the velocities start at zero with $\mathbf P^i=\mathbf X^i$, and let $i^*$ be the particle with the best initial layout, so that $\mathbf X^{i^*}=\mathbf P^{i^*}=\mathbf G$.
(a)~As long as $\mathbf G$ does not change, particle $i^*$ keeps $\mathbf V^{i^*}=\mathbf 0$ and $\mathbf X^{i^*}=\mathbf G$: it re-evaluates $\mathbf G$ once per iteration and evaluates no other point.
(b)~The first move of every other particle is $\mathbf X^i\leftarrow\mathbf X^i+c_2\boldsymbol\rho_2\circ(\mathbf G-\mathbf X^i)$, clipped to the box: turbine $k$ of layout $i$ moves a random fraction in $[0,c_2]=[0,1.50]$ of the way, per coordinate, toward turbine $k$ of $\mathbf G$, although the numbering of the turbines of two independently generated layouts is arbitrary.
\end{proposition}
\begin{proof}
(a)~By induction: with $\mathbf V^{i^*}=\mathbf 0$ and $\mathbf P^{i^*}=\mathbf X^{i^*}=\mathbf G$, all three terms of~\eqref{eq:pso-v} vanish, the particle stays at $\mathbf G$, and $F_p(\mathbf G)$ is no strict improvement, so neither $\mathbf P^{i^*}$ nor $\mathbf G$ changes. (b)~At the first move, $\mathbf V^i=\mathbf 0$ and $\mathbf P^i=\mathbf X^i$, so only the last term of~\eqref{eq:pso-v} remains.
\end{proof}
A DE trial takes coordinate $j$ from the mutant $\mathbf x_{r_1}+F(\mathbf x_{r_2}-\mathbf x_{r_3})$ if $j=j_{\rm rand}$ or $U_j<\mathrm{CR}$, i.e., $1+\mathrm{Bin}(n-1,\mathrm{CR})$ coordinates, on average 27.1 of $n=30$ for $N=15$, again regardless of the turbine numbering. By Lemma~\ref{lem:feas-improve}, $\mathbf G$ moves only if such a combined layout is feasible and better; that this never happened in our runs is an observation, not a consequence of the proposition.
% CHECK-THEORY [T11]: DE takes most coordinates from the mutant (mean 1+CR(n-1) = 27.1 of 30 for N = 15, CR = 0.9; PSO moves all). Lemma lem:feas-improve = Lemma S-lem:S-feas-improve and Proposition prop:feasible-start = Proposition S-prop:S-frozen (a), (b) (optA/supp_theory.tex, S-th-feas); the stall is observed (D13/D14), not implied.
% Page budget: kept in the supplementary material: \SuppTable{tab:D-feasstart}

In Table~\ref{S-tab:D-feasstart}, which logs every candidate, the only feasible PSO candidates (3.33\%) are the re-evaluations of $\mathbf G$ of part~(a), no personal best is updated, and no DE trial is feasible. The numbering of part~(b) is not the whole explanation: after relabeling the turbines of each start to their nearest turbines of $\mathbf G$, the first PSO move is feasible in only 0.1\% of the draws, a convex combination of two layouts in 16.6\%; independent weights per coordinate, with overshoot, place the turbines of two distinct dense layouts (median RMS distance 291~m) too close together (median minimum spacing after the first move 95~m, limit 308~m) or outside the site.
% CHECK-DIAG [D14]: every feasible PSO candidate is the G-particle (zero velocity) re-evaluating G; no personal best updated; DE no feasible trial, none accepted
% CHECK-DIAG [D15]: after optimal relabelling the first PSO move is feasible in < 1 % of draws; only a common weight (convex combination) more often (\NDFsMixScalar %)
% CHECK-DIAG [D16]: feasible initial layouts distinct (matched RMS distance to G > 50 m; median \NDFsInitRms m)

\subsection{Horns Rev~1 16-Turbine Block}
\label{sec:hornsrev-site}
% >>>>> begin analysis/mpce_tab_hr16.tex
%% generated by mpce_results.py -- do not edit by hand
\begin{table}[!t]
\centering
\caption{Horns Rev~1 16-turbine block (5$^\circ$ direction bins; pool of 11 methods: the ten of the budget study and GA): AEP, feasibility, AEP loss and rank.}
\label{tab:hr-site}
\scriptsize\setlength{\tabcolsep}{3pt}
\begin{tabular}{lcccccccc}
\toprule
& \multicolumn{2}{c}{6,030 evaluations, R} & \multicolumn{4}{c}{Loss (\%)} & \multicolumn{2}{c}{Rank, R} \\
\cmidrule(lr){2-3}\cmidrule(lr){4-7}\cmidrule(lr){8-9}
Layout / method & AEP & Feas. & 6k R & 6k F & 30k & 120k & 6k & 30k \\
\midrule
Installed layout & 139.82 & -- & 6.04 & -- & -- & -- & -- & -- \\
\midrule
GA & 138.12 & 30/30 & 7.18 & n/r & \textbf{6.16} & n/r & 2 & \textbf{1} \\
PSO-VNS & \textbf{138.29} & 30/30 & \textbf{7.07} & 6.97 & 6.52 & 6.40 & \textbf{1} & 2 \\
SSA-VNS & 137.59 & 30/30 & 7.54 & 6.98 & 7.19 & 6.47 & 4 & 7 \\
LX-SSA-VNS & 137.41 & 30/30 & 7.66 & 7.02 & 7.07 & 6.44 & 6 & 6 \\
RS-VNS & 137.58 & 28/30 & 7.54$^{28}$ & n/r & 7.04 & 6.25 & 5 & 5 \\
VNS & 137.35 & 30/30 & 7.70 & \textbf{6.91} & 6.85 & 6.23 & 7 & 4 \\
PSO & 137.79 & 19/30 & 7.40$^{19}$ & 8.35 & 6.78$^{28}$ & 6.63 & 3 & 3 \\
SSA & 136.25 & 30/30 & 8.44 & 7.83 & 8.17 & 7.62 & 9 & 9 \\
LX-SSA & 136.00 & 28/30 & 8.61$^{28}$ & 7.89 & 8.50 & 8.28 & 10 & 10 \\
DE & -- & 0/30 & -- & 8.35 & 7.40$^{8}$ & \textbf{6.12} & 11 & 11 \\
MS-SLSQP & 136.31 & 30/30 & 8.40 & 8.00 & 7.82 & 7.31 & 8 & 8 \\
\bottomrule
\end{tabular}
\par\vspace{2pt}\parbox{\columnwidth}{\scriptsize Wake-free AEP 148.81 GWh/yr. AEP: mean (GWh/yr) over the feasible runs of 30 seeds at 6,030 evaluations, random initialization (R); Feas.: feasible runs. Loss: mean AEP loss (\%) of the feasible runs at 6,030 evaluations with random (R) and feasibility-preserving (F) initialization and at 30,030 and 120,030 evaluations (R; 10 seeds at 120,030); superscript: feasible runs when not all are feasible; bold: best method; n/r: not run. Rank: feasibility-aware rank (rule of Table~\ref{tab:friedman68}) in the pool of all 11 methods; seed-paired Holm-adjusted tests in this pool: Table~\ref{S-tab:S-r3-hr-pool}.}
\end{table}
% <<<<< end analysis/mpce_tab_hr16.tex

The model (Section~\ref{S-sec:S-hrmodel}) uses the Jensen wake with a hub-center in-wake test, the thrust coefficient $C_T$ at the free-stream speed, and 5$^\circ$ direction and 1-m/s speed bins; the boundary violation in~\eqref{eq:penalty} is the distance (m) outside the parallelogram. For the installed 80-turbine farm at identical bins, PyWake~2.6.20's NOJ model~\cite{PyWake} gives 661.39~GWh/yr and ours 666.75~GWh/yr (0.81\% more; wake loss 10.39\% (ours) vs.\ 11.11\% (PyWake); Table~\ref{S-tab:F-pywake}); with $C_T$ at the local speed, our model reproduces PyWake's hub-center NOJ (momentum-theory induction) exactly, and rotor averaging does not explain the difference.
% CHECK-DIR [F31] [F32]: PyWake 2.6.20 NOJ(k=0.04) at bins identical to ours (same wake-free AEP): 661.39 vs ours 666.75 GWh/yr (+0.81 %), wake loss 10.39 % (ours) vs 11.11 % (PyWake)
% CHECK-DIR [F33] [F34]: rotor averaging does not explain the difference; with C_T at the local (waked) speed our model reproduces PyWake's NOJ with hub-centre test (momentum-theory induction) to 1e-6 GWh/yr; the old 662.5 constant (\NHRPyWakeDiff) is dropped [F36]
% CHECK-DIR [F35]: installed 16-turbine block: PyWake NOJ 139.08 vs ours 139.82 GWh/yr (+0.53 %)

The installed layout, a regular $7D$ grid, reaches 139.82~GWh/yr (wake loss 6.04\%; Table~\ref{tab:hr-site}). At 6,030 evaluations from random starts, the robust result is feasibility: PSO-VNS is feasible in 30/30 runs, PSO in only 19/30 (exact McNemar $p=\ensuremath {9.8\times 10^{-4}}$; 30/30 with feasible starts) and DE in 0/30. The best layout of all methods at this budget, a PSO-VNS layout (139.78~GWh/yr; Figs.~\ref{S-fig:S-hr16} and~\ref{S-fig:S-layouts-sites}), stays below the installed block, like every random-start run at this budget (GA included).
% CHECK-EXTRA [X52]: HR feasibility PSO-VNS 30/30 vs PSO 19/30, 11 discordant seeds all favour PSO-VNS, exact McNemar p = 9.8e-4
% CHECK-FINAL [C20]: PSO feasible in fewer than 30 of 30 random-start runs ("only"); [C31]: 6,030 PSO-VNS mean AEP below the installed block; no 6,030 random-start run above the installed block with 5-deg bins (Table S-F-hr; 16 runs above at 5 deg over all budgets, 0 at 1 deg: [F21])
% Over the page budget (+0.3 pages in the private build), kept in the supplement (Fig. S-fig:S-layouts-sites):
% \begin{figure}[!htbp]
% \centering
% \pipefig{\textwidth}{figures_mpce/layouts_iea37_hr16.pdf}{5cm}
% \caption{Best layouts. Left and middle: IEA37 Case Study~1, 16 and 36 turbines, best feasible published layout (strict convention) and best PSO-VNS layout at 30,030 evaluations. Right: Horns Rev~1 16-turbine block (5\textdegree{} bins), installed layout ($7D$ grid) and best PSO-VNS layout at 6,030 evaluations (minimum spacing $4D$).}
% \label{fig:layouts-sites}
% \end{figure}
% Over the page budget (+0.35 pages in the private build), kept in the supplement (Fig. S-fig:S-hr16):
% \begin{figure}[!htbp]\centering\pipefig{0.75\textwidth}{figures_mpce/hr16.pdf}{6cm}\caption{...}\label{fig:hr16}\end{figure}

In one pool of the ten methods and GA (Table~\ref{tab:hr-site}; seed-paired tests: Table~\ref{S-tab:S-r3-hr-pool}), PSO-VNS is best at 6,030 evaluations (GA second, $p_{\rm Holm}=0.53$; better than the other nine, $p_{\rm Holm}\le5.1\times10^{-4}$) and GA at 30,030 (139.64~GWh/yr; better than all ten, $p_{\rm Holm}\le1.1\times10^{-4}$; PSO-VNS second, better than the remaining nine, $p_{\rm Holm}\le0.009$). We re-evaluated every final layout with 1$^\circ$ bins and with PyWake's NOJ at the same bins (Tables~\ref{S-tab:F-hr} and~\ref{S-tab:S-r3-hr-eval}). At 6,030 evaluations, GA and PSO-VNS do not differ under any evaluator ($p\ge0.36$), and the AEP advantage of PSO-VNS over PSO on the 19 jointly feasible seeds (+0.54~GWh/yr, $p=0.016$) vanishes with 1$^\circ$ bins (\ensuremath {-}0.05, $p=0.65$) and under PyWake. At 30,030 evaluations, GA's lead holds with 1$^\circ$ bins (+0.48~GWh/yr over PSO-VNS, $p=2.8\times10^{-4}$) and under PyWake (+0.43, $p=4.2\times10^{-4}$); PSO-VNS has the highest mean AEP of the ten other methods under all three evaluations and beats PSO on the 28 jointly feasible seeds also with 1$^\circ$ bins (+0.28~GWh/yr, $p=\ensuremath {6.5\times 10^{-4}}$) and under PyWake; at 120,030 evaluations (10 seeds, GA not run), DE has the highest mean AEP. The finer rose costs the installed block 0.32~GWh/yr and the layouts of the ten methods 1.08~GWh/yr on average, the VNS-based methods most (GA: 0.99 and 1.45~GWh/yr at 6,030 and 30,030 evaluations): like those of the benchmark (Section~\ref{sec:direction}), the optimized layouts exploit the gaps between the coarse direction bins. With 1$^\circ$ bins, none of the 961 feasible optimized layouts (all methods and budgets, GA included) exceeds the installed block (22 with 5$^\circ$ bins), although the looser spacing $4D$ (installed: $7D$) enlarges their feasible set; the comparison tests the optimizers, not the installed design (Section~\ref{sec:disc-practice}).
% CHECK-DIR [F24]: 6,030 PSO-VNS - PSO AEP on 19 jointly feasible seeds: significant only with 5-deg bins; with 1-deg bins and PyWake NOJ PSO has the (n.s.) higher mean
% CHECK-DIR [F26]: 30,030: PSO-VNS highest mean AEP of the ten methods with 5-deg, 1-deg (both phases) and PyWake NOJ; beats PSO on 28 jointly feasible seeds with 1-deg bins and PyWake (p 6.5e-4, 3.8e-4). GA (pool of 11, rev3_sites.json hr_pool/hr_eval, Tables S-r3-hr-pool, S-r3-hr-eval) is first at 30,030 under all three evaluators
% CHECK-DIR [F29]: 120,030: DE highest mean AEP under 5-deg bins, 1-deg bins and PyWake NOJ
% CHECK-DIR [F23]: installed block loses \NFHRInstDrop GWh/yr, optimized layouts \NFHRMeanDrop on average; every VNS-based method loses more than every other method (ten methods; GA drop from rev3_sites.json mean_drop_5to1_gwh)
% CHECK-DIR [F21]: with 1-deg bins 0 of 901 feasible runs of the ten methods exceed the installed block (16 with 5-deg bins); with GA (rev3_sites_reeval.csv): 0 of 961 (22 with 5-deg bins). The former "runs exceed the installed layout" (C30) and the 6,030 run-level test / highest mean AEP (C19, C39) are no longer asserted (R4-1, R4-10; D16, D19).
% Over the page budget (+1.0 and +0.6 pages in the private build), kept in the supplement: \SuppTable{tab:F-hr} \SuppTable{tab:F-pywake}

\subsection{Lillgrund 16-Turbine Block}
\label{sec:lillgrund}
The eight methods of the main comparison and RSD-VNS were also run on a second site with a measured wind climate: a $4\times4$ block of the Lillgrund offshore wind farm (Sweden; Siemens SWT-2.3-93, $D=93$~m), with the 12-sector Weibull climate of PyWake's Lillgrund example~\cite{PyWake}, fitted to seven months of met-mast measurements at 61~m (hub height 65~m)~\cite{Gocmen2016}, the model and bins of Section~\ref{sec:hornsrev-site}, the parallelogram through the corner turbines as boundary, and 30 seeds from random starts at 6,030 and 30,030 evaluations (Section~\ref{S-sec:S-rev-site}, Table~\ref{S-tab:S-rev-site}). The minimum spacing is $3D$, because 16 turbines do not fit into the block at $4D$ (installed minimum spacing: $3.29D$). The feasibility advantage of PSO-VNS over PSO replicates: PSO-VNS is feasible in 17/30 runs at 6,030 evaluations and 30/30 at 30,030, PSO in 0/30 and 19/30 (exact McNemar $p=1.5\times10^{-5}$ and $9.8\times10^{-4}$), and PSO-VNS ranks first at both budgets (mean AEP 113.07 and 114.59~GWh/yr). At 6,030 evaluations, however, only MS-SLSQP is feasible in all 30 runs, and in seed-paired tests it is better than PSO-VNS ($p_{\rm Holm}=0.010$), whereas at 30,030 PSO-VNS is better than every other method except VNS ($p_{\rm Holm}<0.02$; rank and paired test: Section~\ref{sec:archive-reliability}). No method reaches the installed block on average (116.10~GWh/yr; one PSO-VNS layout exceeds it by 0.05\%, also with 1$^\circ$ bins). Re-evaluated with 1$^\circ$ bins (Table~\ref{S-tab:S-r3-lg-eval}), the installed block loses 1.65\% and the optimized layouts 1.14\% on average (PSO-VNS at 30,030: 1.69\%); PSO-VNS keeps the highest mean AEP at 30,030 (112.64~GWh/yr) but is no longer significantly better than VNS or MS-SLSQP ($p_{\rm Holm}=0.58$), and at 6,030 MS-SLSQP has the higher mean (112.10 vs.\ 111.97~GWh/yr). The feasibility results do not depend on the evaluator.

\subsection{IEA37 Case Study~1}
\label{sec:iea37}
% >>>>> begin analysis/mpce_tab_iea37.tex
%% generated by mpce_results.py -- do not edit by hand
\begin{table}[!t]
\centering
\caption{IEA37 Case Study~1: AEP (MWh) of the best run of 11 of the 13 methods of the pool (all 13 with means, ranks and seed-paired tests: Table~\ref{S-tab:S-r3-iea-pool}) and of published layouts~\cite{Baker2019,IEA37repo}.}
\label{tab:iea37}
\scriptsize\setlength{\tabcolsep}{2.5pt}
\begin{tabular}{lcccc}
\toprule
& \multicolumn{2}{c}{16 turbines ($r=1300$~m)} & \multicolumn{2}{c}{36 turbines ($r=2000$~m)} \\
\cmidrule(lr){2-3}\cmidrule(lr){4-5}
Method / source & 6,030 & 30,030 & 6,030 & 30,030 \\
\midrule
Example layout & \multicolumn{2}{c}{366{,}941.6} & \multicolumn{2}{c}{737{,}883.1} \\
Best published, strict$^{a}$ & \multicolumn{2}{c}{418{,}924.4} & \multicolumn{2}{c}{863{,}676.3} \\
Best published, projected$^{b}$ & \multicolumn{2}{c}{421{,}451.1} & \multicolumn{2}{c}{882{,}382.8} \\
\midrule
PSO-VNS & 402{,}782.4 & 413{,}331.9 & 787{,}708.9 & 827{,}385.0 \\
PSO & 403{,}753.1 & 412{,}702.4 & 787{,}988.7 & 830{,}256.4 \\
SSA-VNS & 403{,}251.3 & 406{,}924.6 & 766{,}803.3 & 819{,}930.5 \\
SSA & 395{,}150.2 & 401{,}459.3 & 750{,}707.4 & 802{,}538.2 \\
LX-SSA & 393{,}284.6 & 398{,}770.0 & 746{,}301.3 & 789{,}311.3 \\
DE & 357{,}391.3 & 397{,}402.7 & -- & -- \\
VNS & 400{,}169.3 & 404{,}076.0 & 785{,}791.3 & 824{,}230.2 \\
MS-SLSQP & 405{,}433.8 & 410{,}349.5 & 748{,}571.1 & 833{,}365.0 \\
GA & 399{,}691.3 & 409{,}194.2 & 795{,}530.1 & 836{,}526.4 \\
MS-SLSQP, analytic & 411{,}941.6 & 414{,}406.5 & 839{,}509.3 & 846{,}018.5 \\
PSO-SLSQP, analytic & \textbf{415{,}704.0} & \textbf{418{,}757.4} & \textbf{842{,}857.1} & \textbf{847{,}780.3} \\
\bottomrule
\end{tabular}
\par\vspace{2pt}\parbox{\columnwidth}{\scriptsize AEP of the best feasible run of each method (30 seeds; bold: best method; --: no feasible run; not shown: RS-VNS and LX-SSA-VNS). $^{a}$Best published layout with every turbine within 1~mm of the boundary (participant~4 in both scenarios). $^{b}$Best published layout after projecting the turbines outside the boundary radially onto it, spacing re-checked (participant~12 in both scenarios; up to 3.5~m and 4.9~mm outside before projection).}
\end{table}
% <<<<< end analysis/mpce_tab_iea37.tex

In IEA37 Case Study~1 the participants chose their own optimizers and budgets~\cite{Baker2019}. None of our methods reaches the best feasible published layouts (Table~\ref{tab:iea37}; layouts: Fig.~\ref{S-fig:S-layouts-sites}). Table~\ref{tab:iea37} counts a published layout as feasible at a tolerance of 1~mm (ours: $10^{-6}$~m); after radial projection onto the boundary, 5 more are feasible, and the best are then those of participant~12. Against these, the best PSO-VNS layout at 30,030 evaluations differs in AEP by \ensuremath {-}1.93\% (16 turbines) and \ensuremath {-}6.23\% (36 turbines; strict: \ensuremath {-}1.33\% and \ensuremath {-}4.20\%), in wake loss by +1.73 and +5.21~pp, and would rank fourth among 12 and ninth among 10 feasible submissions. With 36 turbines, the best layout without exact gradients comes from GA (\ensuremath {-}5.20\%), that of the ten original methods from MS-SLSQP (\ensuremath {-}5.56\%). At 6,030 evaluations the gaps are larger (best PSO-VNS layout, projected: \ensuremath {-}4.43\% and \ensuremath {-}10.73\%; all gaps: Table~\ref{S-tab:F-iea-gap}), so they depend more on the budget than on the feasibility convention. As the participants' budgets are not reported comparably, the size of the gap is indicative only; like the published layouts, ours are optimized for 16 directions and one wind speed.
% CHECK-DIR [F41] [F42]: 5 published layouts violate only the boundary (all but one by < 1 cm) and are feasible after radial projection; the best feasible published layout is then participant 12 in both scenarios (421,451.1 and 882,382.8 MWh) instead of participant 4
% CHECK-DIR [F43] [F44]: strict gaps reproduce \NIEA... (-1.33 / -4.20 %); projected gaps of the best PSO-VNS layout at 30,030: -1.93 / -6.23 % (wake loss +1.73 / +5.21 pp); rank fourth of 12+1 and ninth of 10+1; 6,030 projected gaps -4.43 / -10.73 % (Table F-iea-gap)
% verified (manual): budget effect on the best PSO-VNS gap (projected, 6,030 -> 30,030) 2.50 / 4.50 pp vs convention effect (30,030, strict -> projected) 0.60 / 2.03 pp
% CHECK-DIR [F45] + CHECK-FINAL [C48]: 36 turbines, 30,030: best layout of the ten original methods from MS-SLSQP (-5.56 % projected); GA -5.20 % (rev2_summary.json)
% CHECK-FINAL [C40]: best of all methods below the best feasible published layout (strict) in both scenarios
% Over the page budget (+0.3 pages in the private build), kept in the supplement: \SuppTable{tab:F-iea-gap}

\emph{Gradient cost.} With forward differences, every MS-SLSQP gradient costs $2N$ evaluations besides the objective value, $2N+1$ in all, so 6,030 evaluations allow at most 182 gradients with 16 turbines and 82 with 36. Exact or algorithmic-differentiation gradients remove most of this cost, which grows with $N$~\cite{Guirguis2016,Rodrigues2024}, and a smooth wake model, such as the Gaussian wake of IEA37, makes gradients informative~\cite{Thomas2023}. In our runs, MS-SLSQP has, of the ten methods, the best 16-turbine layout at 6,030 evaluations and the best 36-turbine layout at 30,030, but falls behind PSO-VNS with 36 turbines at 6,030 (Table~\ref{tab:iea37}), consistent with this cost. We therefore added SLSQP with the exact analytic gradient of the IEA37 AEP, as multistart MS-SLSQP and after the PSO phase of PSO-VNS (PSO-SLSQP), with each gradient charged three evaluations, its measured cost (30 seeds, random starts; Section~\ref{S-sec:S-rev-grad}, Tables~\ref{S-tab:S-rev-grad} and~\ref{S-tab:S-rev-grad-calls}); both are feasible in all 240 runs. In one pool of all 13 methods (Table~\ref{tab:iea-means}; tests: Table~\ref{S-tab:S-r3-iea-pool}), an exact-gradient method ranks first in every scenario and budget and is better than each of the 11 other methods in all 30 seed pairs ($p_{\rm Holm}=2.2\times10^{-8}$); the two exact-gradient methods differ significantly only at 30,030 evaluations. With 36 turbines, their mean AEP at 6,030 evaluations (832{,}893.3 and 832{,}404.0~MWh) exceeds that of every other method at 30,030 (highest: GA, 823{,}061.5~MWh), and the best PSO-SLSQP layouts at 30,030 lie 0.04\% (16 turbines) and 1.84\% (36 turbines) below the best feasible published layouts (strict). Without exact gradients, PSO-VNS ranks first with 16 turbines (not different from GA at 30,030, $p_{\rm Holm}=0.26$) and GA with 36 turbines, where PSO-VNS ranks fifth and fourth of 13 and does not differ significantly from VNS. The exact-gradient variants differentiate both the objective and the constraints, whereas the original baseline uses finite differences, and PSO-SLSQP also restarts from the incumbent and personal bests: this comparison tests complete implementations of the second phase, not the isolated effect of the AEP derivative. The objective is only piecewise smooth: the derivative treats the upstream-wake indicator as locally constant. Per run, the exact-gradient methods take about 3--6 times as long as PSO-VNS (Section~\ref{sec:runtime}); SLSQP subproblems, constraint evaluations and their Jacobians are not charged, so the comparison is at equal charged budget in objective-equivalent units, not at equal elapsed time or total computation. Their constraints are active at the optimum, so the rounded stored coordinates cannot decide their strict feasibility labels; full-precision reruns do (Section~\ref{sec:archive-reliability}).

\begin{table}[!t]\centering\scriptsize\setlength{\tabcolsep}{1.6pt}
\caption{IEA37 Case Study~1: conditional mean AEP, feasibility and rank in the pool of 13 methods (11 shown).}\label{tab:iea-means}
\begin{tabular}{lrcrrcrrcrrcr}\toprule
& \multicolumn{6}{c}{16 turbines} & \multicolumn{6}{c}{36 turbines} \\
\cmidrule(lr){2-7}\cmidrule(lr){8-13}
& \multicolumn{3}{c}{6,030 evaluations} & \multicolumn{3}{c}{30,030 evaluations} & \multicolumn{3}{c}{6,030 evaluations} & \multicolumn{3}{c}{30,030 evaluations} \\
\cmidrule(lr){2-4}\cmidrule(lr){5-7}\cmidrule(lr){8-10}\cmidrule(lr){11-13}
Method & Mean AEP & Feas. & Rank & Mean AEP & Feas. & Rank & Mean AEP & Feas. & Rank & Mean AEP & Feas. & Rank \\
\midrule
PSO-VNS & 395{,}031.8 & 30/30 & 3 & 402{,}583.1 & 30/30 & 3 & 736{,}832.4 & 30/30 & 5 & 807{,}446.0 & 30/30 & 4 \\
PSO & 390{,}540.2 & 30/30 & 5 & 401{,}347.0 & 30/30 & 5 & 729{,}188.5 & 27/30 & 6 & 802{,}252.3 & 30/30 & 5 \\
GA & 391{,}037.6 & 30/30 & 4 & 401{,}766.1 & 30/30 & 4 & 774{,}971.9 & 30/30 & 3 & 823{,}061.5 & 30/30 & 3 \\
SSA-VNS & 388{,}562.9 & 30/30 & 6 & 396{,}920.5 & 30/30 & 7 & 728{,}870.3 & 30/30 & 7 & 791{,}404.4 & 30/30 & 7 \\
SSA & 378{,}073.1 & 30/30 & 11 & 392{,}875.6 & 30/30 & 11 & 717{,}258.5 & 30/30 & 8 & 771{,}298.0 & 30/30 & 10 \\
LX-SSA & 373{,}977.4 & 30/30 & 12 & 390{,}267.4 & 30/30 & 12 & 709{,}660.5 & 30/30 & 10 & 752{,}030.3 & 30/30 & 11 \\
DE & 338{,}610.6 & 27/30 & 13 & 359{,}952.3 & 30/30 & 13 & -- & 0/30 & 13 & -- & 0/30 & 13 \\
VNS & 387{,}637.3 & 30/30 & 7 & 396{,}192.0 & 30/30 & 8 & 741{,}565.6 & 30/30 & 4 & 801{,}797.4 & 30/30 & 6 \\
MS-SLSQP & 382{,}605.2 & 18/30 & 10 & 394{,}567.0 & 29/30 & 10 & 733{,}342.0 & 4/30 & 12 & 798{,}913.0 & 14/30 & 12 \\
MS-SLSQP, analytic & 407{,}721.0 & 30/30 & 2 & 409{,}976.3 & 30/30 & 2 & \textbf{832{,}893.3} & 30/30 & \textbf{1} & 838{,}852.9 & 30/30 & 2 \\
PSO-SLSQP, analytic & \textbf{407{,}891.4} & 30/30 & \textbf{1} & \textbf{413{,}179.2} & 30/30 & \textbf{1} & 832{,}404.0 & 30/30 & 2 & \textbf{842{,}556.9} & 30/30 & \textbf{1} \\
\bottomrule\end{tabular}
\par\vspace{2pt}\parbox{\linewidth}{\scriptsize Mean AEP in MWh over the feasible runs (bold: best method); Feas.: feasible runs of 30 seeds (strict $10^{-6}$~m labels; Section~\ref{sec:archive-reliability}); --: no feasible run. Analytic variants: analytic objective and constraint derivatives, $c_g=3$. Rank: feasibility-aware rank (rule of Table~\ref{tab:friedman68}) in the pool of the ten methods of the budget study, GA and the two analytic variants; the two methods not shown (RS-VNS, LX-SSA-VNS) and the seed-paired Holm-adjusted tests: Table~\ref{S-tab:S-r3-iea-pool}.}
\end{table}

\begin{table}[!t]
\centering\footnotesize
\caption{Expanded nine-method benchmark pool. Original eight-method results remain a separate analysis.}
\label{tab:ga-main}
\begin{tabular}{lcccc}\toprule
Method & Avg. rank & Best & Feas. (\%) & $p_W$ vs. GA \\\midrule
PSO-VNS & 2.18 & 21 & 100.0 & 0.101 \\
PSO & 2.40 & 17 & 100.0 & 0.696 \\
GA & 2.74 & 16 & 98.6 & -- \\
SSA-VNS & 4.32 & 0 & 99.6 & $9.29\times10^{-10}$ \\
VNS & 5.21 & 3 & 100.0 & $1.07\times10^{-9}$ \\
SSA & 6.34 & 0 & 97.6 & $6.06\times10^{-11}$ \\
MS-SLSQP & 6.85 & 1 & 100.0 & $2.34\times10^{-10}$ \\
LX-SSA & 7.01 & 0 & 96.9 & $6.06\times10^{-11}$ \\
DE & 7.96 & 0 & 71.4 & $1.14\times10^{-10}$ \\
\bottomrule\end{tabular}
\par\vspace{3pt}\parbox{\linewidth}{\footnotesize 68 cases, 30 paired seeds, 6,030 evaluations. Best: cases won outright. $p_W$: Holm-adjusted case-mean Wilcoxon comparison with GA. Qualification and conditional quality follow Section~\ref{sec:setup-stats}; ranks depend on the method pool. Source: Table~\ref{S-tab:S-rev-ga}.}
\end{table}

\subsection{All-Run Reliability and Coordinate Precision}
\label{sec:archive-reliability}
The qualification rule can favour a method with a better mean among fewer successful runs. On Lillgrund at 6,030 evaluations, PSO-VNS has 17/30 feasible runs and a conditional mean AEP of 113.074~GWh, versus 30/30 and 112.865~GWh for MS-SLSQP, which therefore wins the feasibility-first paired run test although its conditional quality rank is lower. We add an all-run paired endpoint: a feasible run beats an infeasible one; among feasible runs the better objective wins (ties within $10^{-9}$); two infeasible runs tie. Over $n$ seed pairs with $W$ wins and $T$ ties, the score is $(W+T/2)/n$ (0.5: no difference), with bootstrap intervals and an exact sign test over the discordant pairs, Holm-corrected over the comparisons of PSO-VNS; unlike the signed-rank test, it does not order infeasible runs by their spacing.
On the 68 benchmark cases (2,040 seed pairs; 95\% case-bootstrap intervals; Table~\ref{S-tab:S-r3-allrun}), PSO-VNS wins 950, ties 348 and loses 742 against PSO (score 0.551, [0.516, 0.587]; $p_{\rm Holm}=4.7\times10^{-7}$): it is better more often than not, by amounts within the margin, consistent with the seed-level interval of Table~\ref{tab:equiv-main}. Against the other six methods of the main comparison it scores 0.73--0.90 (all intervals above 0.5), and ordering them by this score reproduces their rank order (Table~\ref{tab:friedman68}) except that SSA and MS-SLSQP swap. SSA-VNS scores 0.476 against RSD-VNS ([0.452, 0.500]). On Lillgrund (95\% paired-seed bootstrap; Holm over the 16 comparisons with eight methods at two budgets), PSO-VNS scores 0.783 ([0.700, 0.867]) and 0.933 ([0.833, 1.000]) against PSO at 6,030 and 30,030 evaluations ($p_{\rm Holm}=1.68\times10^{-4}$, $1.22\times10^{-5}$), and 0.333 ([0.167, 0.500]; $p_{\rm Holm}=0.395$) and 0.933 against MS-SLSQP; these intervals describe seed variability on this block (all contrasts: Section~\ref{S-sec:S-archive}).

\emph{Coordinate precision.} The stored coordinates are rounded to three (benchmark, IEA37) or two decimals (Horns Rev~1, Lillgrund). Rechecking the feasibility labels of all 56,540 stored records at the $10^{-6}$~m tolerance with the geometry of each site (Table~\ref{S-tab:S-r3-precision}) confirms 50,826 labels with a slack larger than the rounding error and contradicts none; the other 5,714, among them all 240 exact-gradient runs, cannot be decided from the rounded coordinates because a constraint is active at the optimum. Full-precision deterministic reruns of these records: \TBD{C4 rerun: N bit-identical, labels confirmed/changed}. Replaying the stored objectives, computed from the full-precision layouts, from the rounded coordinates changes some values where the hub-center wake test makes the objective jump (Lillgrund: up to 0.225~GWh/yr). All writers now store 17 significant digits.

\subsection{Computational Cost and Published Results}
\label{sec:literature}
\label{sec:runtime}
Per evaluation, the metaheuristics need 0.23--0.45~ms (median over runs) and MS-SLSQP 0.53~ms, 2.3 times as long as PSO (Table~\ref{S-tab:S-r3-time}; means: Table~\ref{S-tab:cost}). In the batch of the stored MS-SLSQP runs, those with $N\ge8$ took 16--22~s (3.4~ms per evaluation) because of thread contention: re-timed single-threaded, they take 4--7~s, and earlier runs of the same method 3.3--4.0~s. The feasibility-preserving initialization, which is not charged, takes a median of 4.8~s (up to 15.5~s) for the 30 layouts of a run on the six largest cases, more than a 6,030-evaluation PSO-VNS run itself (2.8~s). The 36,190 original runs took 101 CPU-hours, run as four parallel processes. In the revision records (medians per run; Table~\ref{S-tab:S-time}), MS-SLSQP with forward differences takes 1.8 times as long as PSO-VNS in the spacing study and 2.2 and 1.6 times on the Lillgrund block at 6,030 and 30,030 evaluations; on IEA37 at 30,030, exact-gradient MS-SLSQP and PSO-SLSQP take 6.5--6.7 and 3.9--4.3 times as long as the GA. Times from separate batches are indicative. Published Kusiak--Song results~\cite{Kusiak2010,Eroglu2012,Eroglu2013,Bansal2017} are discussed in Section~\ref{S-sec:S-literature}; they are not compared numerically because budgets, constraint handling and reporting differ.
% CHECK-FINAL [C46]: cost per evaluation (metaheuristics similar, MS-SLSQP slower than every metaheuristic)
